\chapter{Weak integrals}
\label{chap:integrals}

Integration enters Michal--Bastiani calculus in two places: through the mean value theorem,
which expresses $f(x + h) - f(x)$ as an integral of the differential along the segment, and
through integrals depending on parameters, which are a basic tool to produce smooth maps.
Without a norm there is no Bochner integral, and without completeness Riemann
sums need not converge. The integral used in infinite-dimensional Lie theory is therefore
defined by duality: a vector is the integral of a curve if every continuous linear functional
sees it as such.

Weak integrals defined in this way are unique, by Hahn--Banach, but they need not exist;
existence is the subject of the next chapter (Chapter~\ref{chap:completeness}).
Everything in the present chapter holds in an arbitrary
locally convex space, with the existence of the relevant integrals as an explicit hypothesis.
We first establish the basic properties of the integral and the change of variables formula,
then the first part of the fundamental theorem of calculus, which needs no hypothesis on $F$
at all, and finally the continuity and differentiability of integrals depending on parameters.
Our reference is \cite[\S1.1]{GN26}.

\section{Definition and basic properties}

\begin{definition}[Weak integral]
  \label{def:weak-integral}
  \lean{LocallyConvex.Analysis.HasWeakIntegral}
  \leanok
  Let $\gamma \colon \R \to F$ and $a, b \in \R$. A vector $y \in F$ is a weak integral of
  $\gamma$ over $[a, b]$ if
  \[
    \varphi(y) = \int_a^b \varphi\bigl(\gamma(t)\bigr)\, dt
  \]
  for every continuous linear functional $\varphi$ on $F$.
\end{definition}

\begin{lemma}[Uniqueness]
  \label{lem:weak-integral-unique}
  \lean{LocallyConvex.Analysis.HasWeakIntegral.unique}
  \leanok
  \uses{def:weak-integral}
  Weak integrals are unique.
\end{lemma}

\begin{proof}
  \leanok
  Two weak integrals have the same image under every continuous linear functional, and these
  separate points by Hahn--Banach.
\end{proof}

We write $\int_a^b \gamma$ for the weak integral when it exists.

\begin{remark}
  In Lean, the function that returns the weak integral is total: its value is $0$ when the
  integral does not exist. It is only used where existence is known.
\end{remark}

\begin{lemma}[Algebraic rules]
  \label{lem:weak-integral-algebra}
  \lean{LocallyConvex.Analysis.HasWeakIntegral.add, LocallyConvex.Analysis.HasWeakIntegral.sub, LocallyConvex.Analysis.HasWeakIntegral.neg, LocallyConvex.Analysis.HasWeakIntegral.const_smul, LocallyConvex.Analysis.hasWeakIntegral_const, LocallyConvex.Analysis.HasWeakIntegral.add_adjacent, LocallyConvex.Analysis.HasWeakIntegral.map}
  \leanok
  \uses{def:weak-integral}
  For continuous curves, weak integrals are linear in the integrand and additive in the
  interval: $\int_a^b \gamma + \int_b^c \gamma = \int_a^c \gamma$. The constant curve $v$ has
  integral $(b - a)v$. If $L \colon F \to G$ is continuous and linear, then
  $L\bigl(\int_a^b \gamma\bigr)$ is the weak integral of $L \circ \gamma$.
\end{lemma}

\begin{proof}
  \leanok
  Apply a functional and use the corresponding rules for real integrals. For $L$, it suffices
  to note that $\varphi \circ L$ is a continuous linear functional on $F$ whenever $\varphi$ is
  one on $G$. This is \cite[1.1.4 and 1.1.8]{GN26}.
\end{proof}

The next lemma replaces the estimate $\|\int_a^b \gamma\| \leq (b - a) \sup \|\gamma\|$ of
the normed case. Since there is no norm, the size of a vector is measured by the closed convex
neighborhoods of $0$ that contain it.

\begin{lemma}[Integrals of small curves are small]
  \label{lem:weak-integral-small}
  \lean{LocallyConvex.Analysis.HasWeakIntegral.mem_smul}
  \leanok
  \uses{def:weak-integral}
  Let $a \leq b$ and let $\gamma$ be continuous on $[a, b]$ with weak integral $y$. If
  $W \subseteq F$ is closed, convex, contains $0$ and contains $\gamma([a, b])$, then
  $y \in (b - a + 1)\, W$.
\end{lemma}

\begin{proof}
  \leanok
  If not, Hahn--Banach separates $y$ from the closed convex set $(b - a + 1)W$ by a functional
  $\varphi$. Bounding $\int_a^b \varphi(\gamma)$ by $(b - a)\sup_W \varphi$ gives a
  contradiction. The constant $b - a + 1$ instead of $b - a$ avoids a separate treatment of
  the case $a = b$. Compare \cite[1.1.9]{GN26}.
\end{proof}

\begin{remark}
  When $F$ is normed and $W$ is the closed ball of radius $r$ centered at $0$, the lemma says
  that $\|y\| \leq (b - a + 1)\, r$ whenever $\|\gamma(t)\| \leq r$ on $[a, b]$. Up to the
  harmless constant, this is the classical estimate. For the applications in this document
  (continuity and convergence statements) only the shape of the bound matters, not the
  constant.
\end{remark}

\section{Change of variables}

\begin{proposition}[Change of variables]
  \label{prop:change-of-variables}
  \lean{LocallyConvex.Analysis.hasWeakIntegral_comp_smul_deriv_iff, LocallyConvex.Analysis.hasWeakIntegral_affine_iff}
  \leanok
  \uses{def:weak-integral}
  Let $\phi$ be a real $\Ck{1}$ function on $[a, b]$ and $\gamma$ continuous on
  $\phi([a, b])$. Then $y$ is a weak integral of $t \mapsto \phi'(t)\,\gamma(\phi(t))$ over
  $[a, b]$ if and only if it is a weak integral of $\gamma$ over $[\phi(a), \phi(b)]$.
\end{proposition}

Note that the statement is an equivalence: the existence of either integral implies the
existence of the other. The case used most often is the affine one,
$\phi(t) = a + t(b - a)$ on $[0, 1]$, which shows that $\int_a^b \gamma$ exists if and only if
$\int_0^1 (b - a)\, \gamma(a + t(b - a))\, dt$ does, and that both coincide.

\begin{proof}
  \leanok
  Apply a functional and use the change of variables formula for real integrals.
\end{proof}

\section{The fundamental theorem of calculus, first part}

\begin{theorem}[Fundamental theorem of calculus, I]
  \label{thm:ftc1}
  \lean{LocallyConvex.Analysis.hasWeakIntegral_sub_of_hasDirDerivAt, LocallyConvex.Analysis.weakIntegral_eq_sub_of_hasDirDerivAt}
  \leanok
  \uses{def:weak-integral, def:dir-deriv}
  Let $\gamma$ be differentiable on $[a, b]$, with continuous derivative $\gamma'$. Then
  $\gamma(b) - \gamma(a)$ is the weak integral of $\gamma'$ over $[a, b]$.
\end{theorem}

\begin{proof}
  \leanok
  For every functional $\varphi$, the real function $\varphi \circ \gamma$ has continuous
  derivative $\varphi \circ \gamma'$, and it suffices to apply the real fundamental theorem of
  calculus. This is \cite[1.1.5]{GN26}.
\end{proof}

We remark that no completeness of $F$ is needed: the integral exists because
$\gamma(b) - \gamma(a)$ is a candidate. In other words, derivatives of $\Ck{1}$ curves are
always integrable; the difficulty of the existence problem lies with curves that are not known
to be derivatives.

\begin{corollary}[Mean value theorem]
  \label{cor:ftc-segment}
  \lean{LocallyConvex.Analysis.ContDiffMBOn.hasWeakIntegral_dMB}
  \leanok
  \uses{thm:ftc1, prop:mean-value-weak}
  If $f$ is $\Ck{1}$ on a set containing the segment $[x, x + h]$, then
  $f(x + h) - f(x) = \int_0^1 \dd f(x + th)(h)\, dt$, the weak integral existing in $F$.
\end{corollary}

\begin{proof}
  \leanok
  This is Proposition~\ref{prop:mean-value-weak}, read as a statement about weak integrals.
  See \cite[1.2.6]{GN26}.
\end{proof}

\section{Integrals depending on parameters}

We now let the integrand depend on a parameter $x$ and study the regularity of
$x \mapsto \int_a^b f(x, t)\, dt$. As in the classical case, the arguments rest on the
compactness of $[a, b]$, which turns pointwise closeness into uniform closeness in $t$; the
uniform estimate is then transferred to the integrals by
Lemma~\ref{lem:weak-integral-small}.

\begin{proposition}[Continuity]
  \label{prop:param-continuity}
  \lean{LocallyConvex.Analysis.continuousOn_weakIntegral}
  \leanok
  \uses{lem:weak-integral-small, lem:weak-integral-algebra}
  Let $X$ be a topological space, $U \subseteq X$ and $T \subseteq \R$ open with
  $[a, b] \subseteq T$, and let $f \colon U \times T \to F$ be continuous. If the weak integrals
  $I(x) = \int_a^b f(x, t)\, dt$ exist for $x \in U$, then $I$ is continuous on $U$.
\end{proposition}

\begin{proof}
  \leanok
  By compactness of $[a, b]$ (the tube lemma), for $x'$ close to $x$ the curve
  $f(x', \cdot) - f(x, \cdot)$ takes values in a given closed convex neighborhood $W$ of $0$ on
  $[a, b]$. By Lemma~\ref{lem:weak-integral-small}, $I(x') - I(x) \in (b - a + 1)W$. Since
  the closed convex neighborhoods of $0$ form a basis, this proves the claim. This is
  \cite[1.1.11]{GN26}.
\end{proof}

\begin{theorem}[Differentiation under the integral sign]
  \label{thm:diff-under-integral}
  \lean{LocallyConvex.Analysis.hasDirDerivAt_weakIntegral}
  \leanok
  \uses{prop:param-continuity, lem:mean-value-inclusion}
  Let $f \colon U \times T \to F$ be continuous, with partial derivatives
  $D(x,t)(h) = \dd f\bigl((x,t), (h, 0)\bigr)$ jointly continuous on $(U \times T) \times E$.
  Suppose that the weak integrals $I(x) = \int_a^b f(x, t)\, dt$ exist for $x \in U$, and that
  $J = \int_a^b D(x, t)(h)\, dt$ exists. Then $I$ has derivative $J$ at $x$ in the direction
  $h$.
\end{theorem}

In other words, $\dd I(x)(h) = \int_a^b D(x, t)(h)\, dt$, as in the classical case, provided
that both sides make sense. The existence of the integrals is part of the hypotheses; in a
space with a suitable completeness property it is automatic (Chapter~\ref{chap:completeness}).

\begin{proof}
  \leanok
  For $s \neq 0$, the difference quotient $\frac{1}{s}\bigl(I(x + sh) - I(x)\bigr)$ is the weak
  integral of $t \mapsto \frac{1}{s}\bigl(f(x + sh, t) - f(x, t)\bigr)$. By the mean value
  inclusion and the tube lemma, for small $s$ this curve is uniformly close to
  $t \mapsto D(x, t)(h)$ on $[a, b]$, and Lemma~\ref{lem:weak-integral-small} gives the
  convergence. This is \cite[1.3.15 and 1.3.23]{GN26}.
\end{proof}
